\section{Анализ предметной области и обзор существующих решений}\label{sec1}

В главе рассматриваются ключевые компоненты, на которые опирается декодирование визуальной информации по сигналам мозга: визуальные энкодеры, наборы данных с одновременной регистрацией фМРТ и ЭЭГ, одномодальные методы декодирования, генеративные модели реконструкции, а также методы анализа пространственно-временных характеристик нейросигналов.

\subsection{Визуальные энкодеры для извлечения семантических признаков}\label{sec1:visual}

Визуальный энкодер является одной из ключевых компонент рассматриваемой модели, поскольку от его качества напрямую зависит, насколько полно удается извлечь семантическую информацию из изображения, включая объекты, их позы и контекст сцены. Для многих задач традиционно используются предобученные с учителем энкодеры, среди которых наиболее распространены сверточные архитектуры VGG-19~\cite{simonyan2015deepconvolutionalnetworkslargescale} и ResNet~\cite{he2015deep}, а также Vision Transformer (ViT)~\cite{dosovitskiy2021imageworth16x16words}. Такие энкодеры дают устойчивые визуальные эмбеддинги, и все же, как показывают недавние работы по контрастивному обучению и самообучению, модели вроде CLIP~\cite{radford2021learning} и DINO~\cite{caron2021emergingpropertiesselfsupervisedvision} на базе ViT формируют заметно более выразительные и обобщаемые представления, чем их аналоги, обученные с учителем. Это преимущество устойчиво проявляется в широком круге задач переноса обучения~--- в том числе при декодировании и реконструкции изображений по нейросигналам. Поэтому в роли визуального энкодера здесь выбран именно CLIP~\cite{radford2021learning}: он эффективен, устойчив к доменным сдвигам и обладает богатой семантической структурой латентного пространства.

\subsection{Наборы данных с одновременной регистрацией фМРТ и ЭЭГ}\label{sec1:datasets}

Постановка и решение задачи реконструкции изображений по мозговой активности невозможны без качественных наборов данных, в которых фМРТ и ЭЭГ регистрируются одновременно и под предъявление визуальных стимулов. За последние годы появилось несколько открытых датасетов, включающих совместную регистрацию этих модальностей. Однако большинство доступных наборов либо исключают визуальную стимуляцию, например, регистрируют активность в состоянии покоя или при аудиостимуляции~\cite{Gu2023, Momenian2024}, либо содержат разнесенные во времени сессии фМРТ и ЭЭГ~\cite{Berezutskaya2022}, либо ориентированы на верификацию технической совместимости методов регистрации, а не на исследование стимул-зависимой нейродинамики~\cite{GallegoRudolf2023}. Из доступных на момент написания работы датасетов лишь два обеспечивают одновременную регистрацию при предъявлении визуальных стимулов: датасет Telesford и др.~\cite{Telesford2023} и датасет Wakeman и др.~\cite{Wakeman2015}. Первый содержит записи от 22 участников при просмотре как синтетических стимулов~--- мерцающего шахматного поля и абстрактных сцен в парадигме Inscapes,~--- так и натуралистических~--- коротких видеороликов. Второй включает данные 16 участников при выполнении задачи распознавания статичных изображений лиц. В работе рассматривается датасет~\cite{Telesford2023} в качестве основного, поскольку разнообразие стимулов~--- от контролируемых синтетических до динамических естественных видео~--- обеспечивает более широкое покрытие визуальных паттернов и способствует обучению обобщаемых моделей реконструкции. Именно на его основе формируются обучающие тройки (фМРТ, ЭЭГ, изображение).

\subsection{Одномодальные методы декодирования визуальной информации}\label{sec1:single}

Недавние достижения в декодировании визуальной информации по фМРТ продемонстрировали высокое качество реконструкции за счет интеграции контрастивного обучения и генеративных моделей~\cite{huang2021fmri}. Авторы~\cite{ozcelik2023naturalscenereconstructionfmri} предлагают двухэтапную модель: на первом этапе данные фМРТ кодируются с помощью регрессии, обучаемой индивидуально для каждого участника, после чего реконструируется грубое изображение с помощью VDVAE; на втором~--- изображение уточняется диффузионной моделью, обусловленной как CLIP-эмбеддингами исходного изображения, так и его текстовым описанием, что обеспечивает высокую семантическую согласованность реконструкции. В работе~\cite{chen2023cinematicmindscapeshighqualityvideo} предлагается метод реконструкции видео по временному ряду фМРТ. Используется контрастивное обучение и CLIP-эмбеддинги, аналогично предыдущему исследованию. Генерация видео осуществляется с помощью модифицированной диффузионной модели с временным вниманием и обуславливанием на эмбеддинги фМРТ. В работе~\cite{scotti2023reconstructingmindseyefmritoimage} предлагается диффузионная априорная модель для уточнения фМРТ-эмбеддингов в CLIP-пространстве, устраняющая разрыв между модальностями и повышающая качество декодирования. На основе ЭЭГ также достигнуто сопоставимое качество реконструкции. Основная идея многих работ аналогична: двухэтапная реконструкция~--- сначала контрастивно обучаются эмбеддинги сигналов ЭЭГ к CLIP-эмбеддингам изображений, затем используются предобученная генеративная модель для реконструкции~\cite{Mishra2022, li2024visual, singh2023eeg2imageimagereconstructioneeg}. В качестве кодировщиков применяют различные архитектуры: индивидуальные LSTM для каждого пациента~\cite{singh2023eeg2imageimagereconstructioneeg}, сверточные сети~\cite{Mishra2022}, а также гибридные архитектуры, сочетающие свертки с механизмом внимания~\cite{li2024visual}. Представленные подходы легли в основу проектирования предлагаемого мультимодального энкодера.

\subsection{Генеративные модели для реконструкции изображений}\label{sec1:generative}

Современные генеративные модели изображений сыграли ключевую роль в продвижении задач визуального декодирования. Наиболее успешными оказались латентно-диффузионные архитектуры, которые работают в компактном скрытом пространстве и допускают гибкое обуславливание на внешние эмбеддинги~\cite{esser2021tamingtransformershighresolutionimage, hatamizadeh2024diffitdiffusionvisiontransformers}. Отдельную линию исследований составляют модели, непосредственно обусловленные на CLIP-эмбеддинги изображений и/или текста, что позволяет использовать общее семантическое пространство для управления генерацией~\cite{radford2021learning, crowson2022vqganclipopendomainimage}. В таких подходах CLIP-пространство выступает универсальным интерфейсом между модулями, кодирующими различные модальности, и генератором. В настоящей работе следуется этой парадигме: для реконструкции изображений используется предобученная латентно-диффузионная модель SDXL (англ. Stable Diffusion XL)~\cite{podell2023sdxlimprovinglatentdiffusion}, обусловленная на CLIP-эмбеддинги через адаптер IP-Adapter (англ. Image Prompt Adapter)~\cite{podell2023sdxlimprovinglatentdiffusion, ye2023ipadaptertextcompatibleimage}, что позволяет напрямую переносить информацию из комбинированных эмбеддингов фМРТ--ЭЭГ в пространство генерации.

\subsection{Пространственно-временные характеристики в анализе нейросигналов}\label{sec1:spatiotemporal}

Точное выделение активных областей мозга, или областей интереса (англ. ROI), составляет ключевой этап декодирования фМРТ: значительная доля вокселей соответствует фону и вносит шум, тогда как за обработку конкретных стимулов отвечают локальные зоны~\cite{poldrack2007region, dorin2024forecasting}. Наряду с пространственной локализацией важна временная динамика сигнала, и для ее описания хорошо зарекомендовала себя риманова геометрия. Она представляет нейросигнал ковариационными матрицами и работает с ними как с точками на многообразии симметричных положительно определенных матриц, а не как с векторами в евклидовом пространстве~\cite{barachant2010riemannian, barachant2011multiclass}. Такой подход успешно применялся к ЭЭГ в задачах классификации моторных образов и распознавания состояний~\cite{congedo2017riemannian, yger2016riemannian, samokhina2022classification}, что мотивирует его перенос на данные фМРТ. Классические методы снижения размерности здесь уступают: метод независимых компонент~\cite{hyvarinen2000independent} требует множества сигналов, а метод главных компонент~\cite{Jolliffe2002} игнорирует неевклидову геометрию ковариационных матриц.

Современные архитектуры глубокого обучения, включая рекуррентные нейронные сети, трансформеры и графовые сверточные сети, демонстрируют высокую эффективность при выявлении сложных зависимостей в многомерных временных рядах~\cite{hochreiter1997long, vaswani2017attention, dai2023classification}. Тем не менее, для достижения высокой обобщающей способности этим моделям требуются обширные обучающие выборки. В контексте нейровизуализации объем доступных данных, напротив, существенно ограничен, поскольку запись активности одного испытуемого, как правило, сводится к единственному, протяженному, временному ряду. Данный факт обуславливает необходимость в разработке методов, способных эффективно извлекать пространственно-временные признаки в условиях малого объема выборки.
